\frontmatter
\begin{titlepage}
\centering
\vspace*{1.4cm}
{\sffamily\bfseries\fontsize{36}{40}\selectfont\color{croqink} Croqtile\par}
\vspace{0.1cm}
{\sffamily\bfseries\fontsize{26}{30}\selectfont\color{croqgreen} Programming Guide\par}
\vspace{0.8cm}
{\Large\itshape A guide to the Croqtile language\par}
\vspace{1.1cm}
\rule{0.72\linewidth}{0.6pt}\par
\vspace{0.9cm}
{\large Version \guideversion{} --- \guidereleasedate\par}
\vfill
{\small\color{croqgray} LancerLab\par}
\vspace{1.0cm}
\end{titlepage}
\thispagestyle{empty}
\cleardoublepage

\chapter*{Preface}
\addcontentsline{toc}{chapter}{Preface}

As large models have become widely used, efficient operator implementations
have become a growing part of the development work. Model architectures may
share a few common designs, but the computations within them continue to
vary. Attention mechanisms, sparsity patterns, positional encodings and
quantisation schemes all introduce new implementation requirements. Even
the same operator may need different kernels for different input shapes or
numerical formats. A shared model design therefore leaves substantial
work to be done at the operator level.

That work also grows with the variety of computing platforms. A program may
need to run on several kinds of machine, each with its own memory organisation
and execution model. Within a single family of machines, new instructions and
programming interfaces can require an existing kernel to be reconsidered
and retuned. Operator fusion adds another source of variation: combining
several operations can reduce intermediate data transfers and launch costs,
but each combination brings its own implementation choices. Supporting these
workloads calls for a continuing supply of kernels and the means to develop
them efficiently.

Writing those kernels is demanding because performance depends on decisions
that interact: tile sizes, data layouts, memory staging, pipeline depth and
cooperation between parallel regions. A change in one can require changes
in several others. Higher-level kernel languages have emerged to improve
author productivity by expressing more of this work with less code. Their
central challenge is to simplify programming while keeping the decisions
needed to reach high performance available to the programmer. When an
abstraction hides those decisions, it also limits the implementations that
can be explored.

Croqtile is a language for high-performance heterogeneous computing built
around this challenge, with human productivity as its starting point.
It gives shapes, tile selection, data movement and coordination direct
language forms, so a programmer can describe how a kernel works in one
compact program. Related sizes and indices are derived from shaped data,
reducing the bookkeeping that otherwise accompanies a structural change.
Symbolic dimensions allow the same source to handle different input sizes,
while compile-time checks identify inconsistent shapes and indices before
execution. These abstractions add no runtime overhead. Together, they make
kernels easier to read and revise while leaving performance decisions
explicit.

The same design also supports AI-assisted kernel development and tuning.
Compact source leaves room in an agent's context for profiling results and
previous attempts. Independently expressed operations let the agent change
a pipeline or a parallel arrangement through a small number of edits, and
compiler diagnostics help it locate and correct invalid combinations.
Croqtile is agent-ready in this practical sense: the language supports
repeated cycles of proposing a change, checking it and measuring its effect.
Its AI-native design carries the goal of efficient kernel development into
an automated workflow.

Croqtile is embedded in C++. In a \file{.co} file, ordinary C++ prepares
shaped buffers and calls a \co@__co__@ function; that function describes
the iterations, tiles, data movement and coordination of the parallel work.
This guide introduces those constructs through examples and shows how to
combine them into complete programs.

\section*{Reading guide}

This guide is for developers who want to read and write Croqtile kernels.
It assumes familiarity with C++ functions, pointers and arrays, and with
the shape of a tensor. The examples begin with basic language constructs
and show how to organise parallel computation, select and move tiles, and
coordinate execution. Working through them will help you read a
\file{.co} program and combine these constructs into kernels of your own.

If you are new to Croqtile, begin with Chapter~\ref{ch:getting-started} to
build the compiler and run a first program. Then read
Chapter~\ref{ch:language-basics} in order. Its short examples start with
function declarations and shapes, then introduce iteration, tiling, data
movement and coordination. Look first at the structure of each program,
then use the explanation to understand the role of each statement. Focus
on how the data is divided, moved and used in computation.

Once you are comfortable with the basic syntax, consult
Chapter~\ref{ch:programming-guide} as needed. It groups the language by
function and explains the semantics, constraints and design choices behind
the constructs. Complete kernels appear alongside the features they use,
so you can study a particular topic or return to it while writing a
program. Appendix~\ref{app:quickref} collects the syntax for quick lookup;
the corresponding sections explain the meaning of a construct and the
conditions for using it.

\section*{Conventions}

Croqtile code is set in a monospaced font. Longer examples appear as
numbered listings, collected in the \emph{List of Code Listings}; short
fragments appear inline, for example \co@s32 [4, 8] output;@ or
\co@dma.copy src => shared@. Three kinds of callout box carry the running
commentary: a \textbf{Chapter roadmap} opens each chapter with what it
covers and where to go afterwards, a \textbf{Key point} marks an idea worth
carrying forward, and a \textbf{Common pitfall} records a rule that is easy
to get wrong. Examples use the \texttt{croq::} namespace for host-side
helpers and \texttt{\_\_co\_\_} for the function that describes parallel
work; the two forms appear together in almost every complete program.

\cleardoublepage
\tableofcontents
\cleardoublepage
\lstlistoflistings
\cleardoublepage
